% =====================================================================
% tesi/capitoli/19-caso-distribuzioni.tex
% =====================================================================
% copre: recreate-pokemon-distributions-events/README.md
% copre: recreate-pokemon-distributions-events/STUDIO-01-distribuzioni-gen3-e-ricreazione.md
% verificato-al-commit: 0529162
% ---------------------------------------------------------------------

\chapter{Caso di studio: le distribuzioni di eventi e la loro ricreazione}
\label{cap:distribuzioni}

Il sesto caso di studio è il primo di questo lavoro che porti con sé una data, e la data non è una scelta di pianificazione ma un fatto esterno: il servizio che costituisce l'unico ingresso verso la piattaforma di destinazione cessa di funzionare il 26 febbraio 2027. Tutti gli altri casi sono limitati da ciò che manca, cioè un componente da acquistare, una misura da compiere, una decisione da prendere; questo è limitato da ciò che finisce.

L'oggetto immediato è la ricreazione, su hardware originale e su cartucce possedute, delle distribuzioni di eventi della terza generazione, cioè degli esemplari che nei primi anni Duemila si ottenevano soltanto portando la propria copia del gioco in un luogo fisico. L'oggetto mediato, che è quello che dà senso al primo, è il completamento di una collezione: la totalità delle specie e delle forme alternative dentro il deposito in rete della piattaforma corrente, come impresa di lungo periodo. Il primo è la gamba mancante del secondo, perché una parte delle specie non è mai stata ottenibile per gioco ordinario ed esisteva soltanto dentro un evento a tempo.

\section{Perché una distribuzione è un problema di conservazione}

Una distribuzione di evento era il canale attraverso il quale un esemplare non ottenibile per gioco ordinario entrava in una cartuccia. Oggi la medesima operazione avviene per rete oppure con un codice alfanumerico; allora richiedeva la presenza fisica, un collegamento fra il proprio apparecchio e quello dell'evento, e una finestra temporale. Gli esemplari così distribuiti erano speciali in un senso tecnico e non decorativo: appartenevano a specie assenti dal gioco, oppure portavano una mossa che quella specie non può apprendere in alcun altro modo, il che li rende non riproducibili nemmeno per gioco esaustivo.

Da qui la natura del problema. Passata la finestra, l'esemplare non è più ottenibile e non lo diventa mai più, perché nessun secondo canale è stato aperto. Ne segue che la ricreazione del meccanismo di distribuzione non è una comodità ma la sola via esistente, e che il suo oggetto non è l'esemplare ma la procedura che lo generava: far ripetere al gioco esattamente ciò che l'apparecchio dell'evento gli faceva fare.

Questo capitolo tratta la sola terza generazione, per tre ragioni concomitanti. È l'ultima priva di un servizio di rete, dunque quella in cui la distribuzione fisica era l'unico canale. È la prima con un formato dati cifrato, dunque quella in cui la costruzione a posteriori di un esemplare valido richiede il macchinario descritto nel capitolo~\ref{cap:cifratura}. Ed è la generazione delle cartucce su cui questo lavoro opera.

\section{Il canale, e ciò che il canale consente}

Una distribuzione di terza generazione viaggia per \term{multiboot}, il meccanismo con cui una console riceve un programma dal collegamento seriale e lo esegue in memoria di lavoro invece che dalla cartuccia, descritto nel capitolo~\ref{cap:multiboot}. Il fatto rilevante per questo capitolo non è il meccanismo, che il lavoro documenta già come canale del ponte fra generazioni, ma la circostanza che le distribuzioni ufficiali impiegassero esattamente quello: il canale che questo progetto studia come via di trasferimento è il canale che il produttore impiegava come via di distribuzione.

La proprietà che rende il multiboot adatto allo scopo è che il programma ricevuto conserva l'accesso al salvataggio della cartuccia inserita, e lo impiega: legge il salvataggio, vi scrive l'esemplare, e conclude la procedura \cite{video-distribuzioni}. Non si tratta dunque di una ricezione per via di gioco, ma di una scrittura esterna dentro il salvataggio, con l'unica differenza, rispetto a una manomissione, che il programma è ufficiale e conosce le regole del formato. La distinzione è esattamente quella che il capitolo~\ref{cap:conversione} enuncia a proposito della coerenza per costruzione contro la coerenza rispetto ai controlli conosciuti.

Due dettagli operativi meritano registrazione, perché appartengono al genere di condizione la cui violazione produce un fallimento silenzioso. L'orientamento del cavo non è indifferente, e il capo di dimensione minore va inserito nell'apparecchio che porta la ROM di distribuzione. Il ricevente va avviato in modalità multiboot tenendo premuti due tasti all'accensione, condizione riconoscibile dalla scomparsa del logo di avvio.

\section{Che cosa contiene una ROM di distribuzione}

Questa sezione descrive un formato che nessun riferimento enciclopedico documenta e che è stato ricostruito per analisi del codice \cite{video-distribuzioni}. Vale esporla con ordine, perché il metodo impiegato è trasferibile e alcune sue conclusioni riguardano oggetti già trattati altrove in questo lavoro.

Il programma da inviare non si trova nella ROM in forma leggibile. Lo si individua confrontando la sequenza di byte che transita sul collegamento con il contenuto della ROM, e ciò che si estrae in tal modo è poco codice, in massima parte dedicato a confermare al mittente la riuscita del trasferimento. In coda a quel codice si trova però una chiamata di interruzione software al sistema, il cui parametro seleziona la funzione richiesta, e le prime parole del blocco che segue dichiarano il metodo di compressione impiegato e la dimensione dei dati una volta espansi: è quanto basta per decomprimerli fuori dal gioco e ottenere il codice effettivo della distribuzione.

\begin{daverificare}
La fonte indica il parametro della chiamata di sistema con il valore 11 e non dichiara la notazione impiegata, dunque non è determinato se si tratti del valore decimale o esadecimale. L'identificazione con la funzione di decompressione documentata dal riferimento dell'hardware richiede il confronto con la tabella delle chiamate di sistema, che non è stato compiuto.
\end{daverificare}

La difesa che il programma oppone alla modifica è un checksum additivo esteso a tutti i byte della ROM, confrontato con un valore precalcolato contenuto nel programma stesso; se il confronto fallisce, il gioco espone un codice di errore e non riceve nulla. Il modo in cui quel controllo è stato individuato merita di essere riportato perché costituisce una tecnica e non un aneddoto: i testi dei messaggi di errore risiedono nella ROM in forma leggibile, e risalendo il puntatore al testo del messaggio si raggiunge il codice che lo produce, e da quello il controllo che lo invoca.

Il legame con il capitolo~\ref{cap:analisi} è diretto e la conclusione va enunciata, perché rovescia l'impressione immediata. Un checksum additivo è invariante rispetto all'ordine degli addendi, e da quella invarianza il presente lavoro deduce che una permutazione dei blocchi lo attraversa con probabilità unitaria. Qui la medesima proprietà ha una conseguenza pratica: un checksum additivo si corregge per differenza, senza necessità di ricalcolarlo, e non oppone resistenza effettiva a una modifica deliberata. Difendeva dalla corruzione del trasferimento, non da un avversario, che è la distinzione stabilita nel capitolo~\ref{cap:integrita}.

I dati dell'esemplare risiedono in una struttura di dimensione contenuta, che porta la specie, il livello, le mosse e un puntatore separato al nome. Il puntatore separato è la ragione per cui è possibile aggiungere mosse senza distruggere il nome, che è invece l'effetto di una scrittura oltre lo spazio riservato.

La scoperta di maggiore portata è però un'altra, ed è che il programma di distribuzione riusa il codice del gioco \cite{video-distribuzioni}. In particolare riusa il meccanismo con cui il gioco imposta i parametri di un esemplare, dove ciascun parametro è identificato da un indice numerico. Ne seguono due conseguenze, e la seconda è quella che conta: si può individuare il punto in cui un parametro viene impostato cercandone l'indice, e si può impostare un parametro diverso sostituendo l'indice. La ricreazione di un esemplare che deve portare un riconoscimento particolare è stata ottenuta riusando il codice che imposta il vincolo di obbedienza e sostituendovi l'indice del riconoscimento. L'insieme dei parametri impostabili non è dunque limitato da quello che il programma originale impostava.

\section{I nove problemi della ricreazione}

I problemi risolti dalla ricreazione costituiscono, presi insieme, la mappa di ciò che occorre conoscere per costruire un evento arbitrario, e ciascuno di essi corrisponde a un campo o a un meccanismo che i capitoli precedenti hanno descritto dal lato della lettura.

Il vincolo di obbedienza è il primo, e va attivo per una specie e spento per tutte le altre; non compariva nella struttura individuata, ed è stato raggiunto per la via degli indici. L'esperienza è il secondo: non è un valore libero ma una funzione del livello e del gruppo di crescita, e il programma porta con sé la tabella dell'esperienza per livello di ciascuno dei sei gruppi, cosicché è sufficiente dichiarare il gruppo corretto. L'abilità è il terzo, e la prima ricreazione ne assegnava una sola, rendendo illegittima la metà degli esemplari prodotti per le specie che ne ammettono due; esiste un contrassegno che, abilitato, fa derivare la scelta dal primo bit del valore di personalità, che è la medesima regola descritta nel capitolo~\ref{cap:identita} vista dal lato di chi genera.

Il sesso dell'allenatore di provenienza è il quarto: normalmente sorteggiato, va reso deterministico per gli esemplari la cui provenienza dichiarata è fissa. La circostanza che fra le distribuzioni note esistano almeno tre formule diverse per quel sorteggio è già un indizio che il codice delle distribuzioni non fosse unico. La lucentezza è il quinto, e mostra la struttura del problema con particolare chiarezza: le distribuzioni ordinarie contengono un ciclo che rigenera il valore di personalità fino a ottenere un esemplare non lucente, dunque la ricreazione di un esemplare che poteva essere lucente richiede di disattivare quel ciclo e accettare il primo valore, mentre la ricreazione di un esemplare sempre lucente richiede di imporre i sedici bit inferiori del valore di personalità.

L'oggetto tenuto è il sesto e misura il grado di fedeltà raggiunto. In un caso l'esemplare porta con probabilità uguale una di due bacche, e la scelta nell'originale non è un sorteggio qualsiasi ma la sesta estrazione del generatore pseudocasuale, divisa per tre, di cui si considera il primo bit; replicarla ha richiesto codice aggiuntivo, e lo spazio è stato ottenuto rimuovendo i messaggi di errore delle lingue non impiegate.

La posta è il settimo ed è il più difficile fra quelli risolti, con oltre cento tentativi dichiarati. Un esemplare che porta una lettera ha due dati in due luoghi distinti: l'oggetto nella struttura dell'esemplare, il contenuto della lettera in una sezione propria del salvataggio, e i due sono associati da un indice che risiede anch'esso nella struttura. Scrivere il contenuto richiede dunque di leggere e riscrivere il salvataggio, non di comporre una struttura, e la procedura descritta è la seguente: individuare la sezione corretta nella copia più recente del salvataggio impiegando l'indice di salvataggio, farne una copia di lavoro, individuare una posizione libera fra quelle riservate agli esemplari in squadra, scrivere il contenuto, ricalcolare il checksum del settore e riscriverlo.

Il numero delle sezioni, la loro rotazione fra le posizioni e l'esistenza di due copie alternate coincidono con quanto il capitolo~\ref{cap:integrita} stabilisce sul sorgente, e la coincidenza vale come conferma indipendente, giacché ottenuta per analisi del codice della ROM di distribuzione e non per lettura del disassemblato.

\begin{daverificare}
La collocazione della posta nella quarta sezione del salvataggio proviene da una fonte di quarto livello e non è stata confrontata con il sorgente, dove il numero delle sezioni e la loro rotazione sono invece verificati.
\end{daverificare}

Gli esemplari distribuiti come uova sono l'ottavo problema, e introducono una struttura di selezione: a ciascuna specie è associato un peso, i pesi hanno somma nota, si estrae un valore nell'intervallo e si sottraggono i pesi in sequenza fino all'esaurimento, e la specie sulla quale la sottrazione si arresta è quella distribuita. Poiché un'uovo acquisisce l'identificativo dell'allenatore al momento della schiusa, il programma deve leggere quell'identificativo dal salvataggio. Una variante lucente impiega la quarta e la quinta estrazione del generatore invece della terza e della quarta, asimmetria priva di spiegazione nota che la ricreazione conserva per fedeltà.

Il seme derivato dall'orologio è il nono e produce una misura degna di nota. Un esemplare era generato da un seme costruito sommando secondi, minuti e ore dell'orologio interno, e da tale costruzione segue che gli esemplari distinti ottenibili erano soltanto duecentoquattordici, cioè un insieme di gran lunga inferiore a quanto la dimensione del seme suggerirebbe: è il medesimo fenomeno di riduzione dello spazio effettivo che il capitolo~\ref{cap:analisi} misura sulla chiave di cifratura. La ricreazione abbandona l'orologio e ricava un seme valido riducendo per modulo quello generato ordinariamente, con l'effetto collaterale di funzionare anche sulle cartucce la cui batteria dell'orologio è esaurita, che è la condizione tipica di un supporto di vent'anni.

\section{I casi non chiusi, e il modo in cui sono trattati}

Tre casi non sono stati determinati, e il modo in cui la comunità li tratta è la parte metodologicamente più istruttiva della fonte, perché coincide con la disciplina che il capitolo~\ref{cap:perimetro} impone a questo lavoro.

Il primo è stato chiuso per ricerca esaustiva. Un insieme di esemplari era stato generato in lotti a partire da semi di origine a sedici bit, e uno di essi non era rigenerabile perché il suo seme non era fra quelli preservati; poiché la capienza del gioco non è multiplo della dimensione del lotto, alcuni esemplari furono liberati per esaurire lo spazio, e con essi il loro seme. La determinazione è avvenuta percorrendo l'intero spazio dei $65\,536$ semi possibili e verificando che uno solo genera l'esemplare osservato. La struttura dell'argomento è quella che il capitolo~\ref{cap:analisi} impiega sulla tabella delle permutazioni: uno spazio abbastanza piccolo da essere percorso per intero, e l'unicità della soluzione che trasforma un'ipotesi in una determinazione.

Il secondo è opposto e resta aperto. Per una distribuzione, nessuna delle formule note per il sorteggio del sesso dell'allenatore si accorda con i campioni conservati, e le ipotesi in campo sono tre: che il valore fosse copiato dal salvataggio del ricevente, che alcuni campioni siano contraffatti, o che quella distribuzione impiegasse una quarta formula. Una ricerca su tutte le combinazioni di scorrimento e divisione ha prodotto una formula compatibile con tutti i campioni, e l'autore della ricerca dichiara di ritenerla improbabile.

\begin{daverificare}
La formula individuata per quel sorteggio, cioè una divisione per il valore \hx{CC0} di cui si considera il primo bit, si accorda con i campioni disponibili ma è dichiarata improbabile dalla fonte stessa. Non va promossa a fatto, e la sua eventuale conferma richiede campioni ulteriori.
\end{daverificare}

Il terzo riguarda esemplari catturati in ambiente di gioco e replicati per un evento di scambio, dei quali tre non sono conservati. I valori reperibili sono compatibili con la generazione soltanto ammettendo una condizione particolare in chi li catturò, che si sarebbe verificata per quei tre e non per gli altri nove, e un campione conservato porta un'abilità che tale condizione escluderebbe. La fonte include provvisoriamente i dati reperiti dichiarandone la provvisorietà e sollecita pubblicamente campioni migliori, che è la forma corretta di trattare un dato non determinato.

\section{L'apparecchio del negozio, e i limiti della replicabilità}

Una parte del corpus proviene da un apparecchio installato in un negozio, aperto nel novembre 2001 e chiuso all'inizio del 2005, le cui campagne di distribuzione si estesero dal novembre 2001 al marzo 2003 per la seconda generazione e dall'agosto 2003 al dicembre 2004 per la terza \cite{video-pcny}. La scarsità della documentazione su di esso ha una causa dichiarata, cioè il divieto di fotografie in vigore nel negozio, e ciò rende la preservazione dell'apparecchio un evento di rilievo documentale oltre che tecnico.

Il contenuto preservato comprende due dischi in un formato proprietario destinato allo sviluppo, leggibile soltanto da lettori dedicati, alcune schede di memoria che dichiarano al programma quali esemplari distribuire e in quale finestra temporale, e schede ulteriori che operano come chiave e in assenza delle quali il programma non procede. Due proprietà del funzionamento vanno registrate perché determinano la replicabilità. L'orologio interno dell'apparecchio deve indicare una data compresa nella finestra della campagna, dunque la finestra temporale era un controllo e non una convenzione. E il trasferimento avviene attraverso un dispositivo di scrittura dedicato, mentre l'apparecchio portatile collegato al cavo ha funzione dimostrativa: la fonte dichiara di aver verificato che con un apparecchio ordinario il trasferimento non avviene.

Ne segue la conclusione che interessa questo capitolo, ed è negativa nel modo utile: quella via di distribuzione non è replicabile senza l'hardware originale, e la sola strada praticabile per gli eventi di quel negozio è la ricreazione per multiboot. Restano documentate due proprietà degli esemplari, cioè la denominazione dell'allenatore di provenienza comprensiva dell'indicazione della postazione, e l'identificativo dell'allenatore incrementato a ogni distribuzione, che la ricreazione ha replicato leggendo il contatore e iniettandone il valore nel programma prima del trasferimento.

Un dato marginale riguarda la conservazione e vale per analogia: i dischi impiegati sono supporti scrivibili e non stampati, dunque soggetti a degrado più rapido, e la fonte raccomanda di duplicarli. È l'argomento che questo lavoro applica alle batterie delle cartucce nel capitolo~\ref{cap:supporto}.

\section{Le vie di iniezione su cartuccia originale}

La parte operativa del caso di studio è la scelta della via con cui un evento entra in una cartuccia posseduta, e le vie sono quattro, selezionabili in funzione dell'hardware disponibile \cite{video-ereader-eventi}. Una premessa le precede, perché riguarda due funzioni del gioco che si somigliano e non coincidono: la ricezione di una carta di dono e la ricezione di un programma di evento sono meccanismi distinti, la prima si sblocca dentro il gioco digitando una frase concordata, e nel titolo in cui la seconda è stata rimossa dai menu essa si riapre soltanto ricevendo, attraverso la prima, un programma che la ripristina.

La prima via impiega l'emulazione e serve a chi possiede la cartuccia e un modo per estrarne il salvataggio, ma non una seconda console. Due condizioni sono dichiarate necessarie e non sono formali: l'emulatore deve impiegare il programma di avvio estratto dall'hardware reale, perché la sua ricostruzione non riproduce tutte le funzioni e gli eventi terminano in errore, e l'orologio interno va disattivato per non alterare la data del salvataggio. La seconda via impiega il dispositivo di lettura delle carte con un salvataggio precostituito di dimensione maggiore, la terza una scheda riprogrammabile con un salvataggio di dimensione minore, e la quarta la stampa delle carte con il codice ottico, dichiaratamente capricciosa fino al punto di richiedere l'impiego deliberato del driver di stampa errato.

Esiste infine la via che non impiega alcun dispositivo di lettura e opera direttamente sul salvataggio estratto, iniettandovi la carta con un editor dedicato. Questa via porta l'avvertimento che il presente lavoro adotta come norma di verifica del caso di studio, e che è registrato nel capitolo~\ref{cap:collaudo}: se il salvataggio contiene già una carta, essa va esportata e conservata prima di qualunque scrittura, perché può trattarsi di un evento che la comunità non ha ancora preservato, e in tal caso la sovrascrittura distrugge un dato unico invece di un dato ricostruibile.

Sull'hardware di estrazione la fonte elenca le vie disponibili e cita espressamente il lettore che questo progetto ha acquistato per il caso di studio del capitolo~\ref{cap:smeraldo}, il che costituisce una conferma indipendente di quella scelta.

Una via ulteriore, che nessuna delle altre fonti nomina, passa dal lettore di cartucce portatili collegato alla console domestica dell'epoca \cite{video-gbi-bonus}. L'interfaccia alternativa scritta dalla comunità per quel dispositivo invia un programma multiboot senza alcun cavo, poiché il collegamento è interno all'apparecchio, ed è in grado di estrarre programma di avvio, ROM e salvataggio e di ripristinare quest'ultimo: l'intero ciclo di modifica avviene dunque sulla console, senza calcolatore e senza lettore esterno. Il contributo di quella fonte all'obiettivo di collezione non è però un esemplare ma un oggetto, giacché il programma originale del disco supplementare consegnava, al superamento di una prova, i titoli di accesso ad alcuni eventi e, su un salvataggio particolare, la carta nautica che conduce all'incontro con una specie altrimenti assente: oggetti irripetibili quanto gli esemplari, e che sbloccano l'unica via legittima per ottenere quelle specie in quei titoli.

\section{La catena verso la piattaforma corrente, e la data}

Questa sezione contiene il vincolo che governa il caso di studio, e le sue asserzioni sono verificate sulla comunicazione ufficiale e sul riferimento enciclopedico.

Il servizio di deposito per la console portatile a doppio schermo cessa il 26 febbraio 2027 alle 12:00 del fuso giapponese, e nel medesimo istante cessa la funzione che trasferisce verso il deposito della piattaforma corrente; il produttore invita a spostare in anticipo quanto si desidera conservare e non annuncia alcun periodo di tolleranza \cite{bank-fine-servizio}. Quel servizio e il suo programma di trasferimento erano rimasti gli ultimi con funzioni in rete su quella console dopo la chiusura generale dell'aprile 2024.

Il fatto che rende la data operativa per questo caso di studio è che il programma di trasferimento accetta come sorgente soltanto i titoli della quinta generazione e le riedizioni virtuali della prima e della seconda \cite{bulbapedia-trasferimenti}. La terza generazione non vi accede, e un esemplare nato su cartuccia deve dunque attraversare quattro passaggi, ciascuno irreversibile: dalla terza alla quarta generazione attraverso la struttura di migrazione, che richiede una console dotata dello slot per le cartucce della piattaforma precedente e impone sei esemplari per volta, la medesima lingua ai due capi, un intervallo minimo fra le migrazioni nei primi titoli e il rifiuto degli esemplari che conoscono una mossa di servizio; dalla quarta alla quinta attraverso il trasferimento senza fili fra due apparecchi, che richiede il completamento della vicenda principale e l'inventario nazionale sul titolo di destinazione, rifiuta le uova e restituisce al mittente gli oggetti tenuti; poi il programma di trasferimento verso il deposito, e infine il deposito verso la piattaforma corrente.

Ne discendono tre conseguenze operative. La prima è un requisito di hardware che il presente lavoro non aveva dichiarato: il primo passaggio richiede una console della generazione dotata dello slot per le cartucce della piattaforma precedente, che la console impiegata nel caso di studio del capitolo~\ref{cap:3ds} non possiede. La disponibilità di tale apparecchio è stata accertata, e il requisito che resta da verificare è di natura diversa e meno visibile, cioè la corrispondenza di lingua fra i titoli ai due capi di ciascun passaggio, la cui violazione non produce un errore ma l'impossibilità di avviare il trasferimento. La seconda è che il secondo passaggio richiede due apparecchi accesi simultaneamente, uno dei quali può essere la console già disponibile. La terza è favorevole, e riguarda i titoli: fra le cartucce del caso di studio del capitolo~\ref{cap:3ds} quattro appartengono alla quarta generazione e una alla quinta, dunque la catena non richiede acquisti di titoli e richiede soltanto una verifica di lingua.

\begin{daverificare}
Fonti secondarie affermano che il completamento integrale della collezione diverrebbe impossibile dopo la cessazione del servizio, per una specie della terza generazione i cui contrassegni non sarebbero riproducibili dalle vie moderne. Se l'affermazione regge, il traguardo raggiungibile dopo quella data è inferiore alla totalità, e il senso della scadenza passa da urgenza a irreversibilità. L'asserzione non è verificata su fonte di secondo livello.
\end{daverificare}

\section{Legittimità, perimetro, e ciò che resta aperto}

Sulla legittimità degli esemplari ricreati la posizione di questo lavoro è la medesima adottata nel capitolo~\ref{cap:automazione} per gli esemplari ottenuti automatizzando il gioco, e va enunciata senza attenuazioni. Una ricreazione fedele riproduce il metodo di generazione dell'originale, dunque produce un esemplare coerente rispetto ai controlli che un verificatore sa compiere, ma non produce l'esemplare distribuito allora: produce una copia costruita con la medesima procedura. La distinzione va dichiarata quando un esemplare viene scambiato o esposto, e va conservata scritta accanto al dato anziché nella memoria di chi lo ha prodotto.

\begin{daverificare}
Non è verificato se un verificatore di legittimità indipendente accetti un esemplare prodotto da una ricreazione fedele, né quali controlli applichi agli esemplari di provenienza da evento. È la questione che decide se il caso di studio raggiunga il proprio obiettivo, e si risolve senza hardware costruendo un esemplare con il pacchetto descritto nel capitolo~\ref{cap:ponte} e sottoponendolo allo strumento.
\end{daverificare}

Restano infine due questioni di perimetro, che questo lavoro registra come aperte anziché risolverle, secondo il criterio del capitolo~\ref{cap:perimetro}. La prima è che l'ultimo tratto della catena impiega due programmi sui quali l'assistenza tecnica di questo progetto è dichiaratamente esclusa, e nessuna via tecnica alternativa esiste, poiché il programma di trasferimento è l'unico ingresso per tutto ciò che precede l'ottava generazione: la contraddizione fra la norma e l'obiettivo è reale e la sua risoluzione non è tecnica. La seconda è che tre delle quattro vie di iniezione richiedono materiale di terze parti, cioè una ROM di distribuzione oppure un salvataggio precostituito, mentre la norma esclude i salvataggi acquisiti in rete; la norma non si applica automaticamente, giacché non si tratta di importare il salvataggio di un titolo della serie, e proprio per questo la decisione va presa esplicitamente e non dedotta.

Fino a quelle decisioni il caso di studio opera su ciò che le precede, che è la parte maggiore del lavoro: la lettura delle fonti, la verifica della legittimità, la determinazione dei dati non chiusi, e la preparazione della procedura di scrittura secondo il protocollo del capitolo~\ref{cap:collaudo}. Il vincolo temporale, per parte sua, non attende alcuna decisione.
